
\begin{frame}{}
\begin{center}
    {\Huge \textbf{Extra slides}}
\end{center}
\end{frame}





% %%%%%%%%%%%%%%%%%%%%%%%%%%%%%%%%%%%%%%%%%%%%%%%%%%%%%%%%%%%%%%%%%%%%%%%
% %%%%%%%%%%%%%%%%%%%%%%%%%%%%%%%%%%%%%%%%%%%%%%%%%%%%%%%%%%%%%%%%%%%%%%%
% %%%%%%%%%%%%%%%%%%%%%%%%%%%%%%%%%%%%%%%%%%%%%%%%%%%%%%%%%%%%%%%%%%%%%%%










%%%%%%%%%%%%%%%%%%%%%%%%%%%%%%%%%%%%%%%%%%%%%%%%%%%%%%%%%%%%%%%%%%%%%%%%%
%%%%%%%%%%%%%%%%%%%%%%%%%%%%%%%%%%%%%%%%%%%%%%%%%%%%%%%%%%%%%%%%%%%%%%%%%
%%%%%%%%%%%%%%%%%%%%%%%%%%%%%%%%%%%%%%%%%%%%%%%%%%%%%%%%%%%%%%%%%%%%%%%%%

\begin{frame}{Selected experimental results.}

% \begin{itemize}
% \item The ``Refit estimate'' column shows the result of re-fitting the model removing
% the points with the largest influence scores.
%
% \item Stars indicate significance at the 5\% level.
%
% \item Refits that achieved the desired change are bolded.
% \end{itemize}

{
\MicrocreditProfitResultsTable{}
}

\onslide<2->{
\CashTransfersResultsTable{}
}

\onslide<3->{
\OHIEResultsTable{}
}
%
A $*$ indicates statistical significance at the 95\% level.\\
%
\onslide<4->{
    In the wild!  \parencite{
    eubank_enfranchisement_2022,
    mart_how_2022,
    finger_adoption_2022,
    turnbull-dugarte_mobilising_2022,
    shiffman:2023:droppingsinglecell,
    huang:2026:chatbotamip}
}

\end{frame}



% %%%%%%%%%%%%%%%%%%%%%%%%%%%%%%%%%%%%%%%%%%%%%%%%%%%%%%%%%%%%%%%%%%%%%%%
% %%%%%%%%%%%%%%%%%%%%%%%%%%%%%%%%%%%%%%%%%%%%%%%%%%%%%%%%%%%%%%%%%%%%%%%
% %%%%%%%%%%%%%%%%%%%%%%%%%%%%%%%%%%%%%%%%%%%%%%%%%%%%%%%%%%%%%%%%%%%%%%%

\input{determination}



% %%%%%%%%%%%%%%%%%%%%%%%%%%%%%%%%%%%%%%%%%%%%%%%%%%%%%%%%%%%%%%%%%%%%%%%
% %%%%%%%%%%%%%%%%%%%%%%%%%%%%%%%%%%%%%%%%%%%%%%%%%%%%%%%%%%%%%%%%%%%%%%%
% %%%%%%%%%%%%%%%%%%%%%%%%%%%%%%%%%%%%%%%%%%%%%%%%%%%%%%%%%%%%%%%%%%%%%%%

\begin{frame}[fragile]{How to compute the influence scores?}

How can we compute the influence scores
%
%\begin{align*}
%
$
\infl_n =
\fracat{\partial \thetafun(\thetahat(\w))}
       {\partial\w_n}{\w = \onevec}
%
$?
%\end{align*}
%

\pause
\vspace{1em}
By the \textbf{chain rule},
$
\infl_n =
\fracat{\partial \thetafun(\theta)}
      {\partial\theta}{\thetahat(\onevec)}
\fracat{\partial \thetahat(\w)}
    {\partial\w_n}{\w = \onevec}
%
$.

\pause
\vspace{1em}
Recall that
$\sumn \w_n G(\thetahat(\w), \d_{n}) =  \zP$ for all $\w$ near $\onevec$.

\vspace{1em}
$\Rightarrow$
By the \textbf{implicit function theorem}, we can write
$\fracat{\partial\thetahat(\w)}{\partial\w_n}{\w=\onevec}$ as a linear system
involving $G(\cdot, \cdot)$ and its derivatives.

\pause
\vspace{1em}
$\Rightarrow$
The $\infl_n$ are automatically computable from $\thetahat(\onevec)$ and
software implementations of $G(\cdot, \cdot)$ and $\thetafun(\cdot)$
using \textbf{automatic differentiation}.

\begin{lstlisting}
> import jax
> import jax.numpy as np
> def phi(theta):
>     ... computations using np and theta ...
>     return value
>
> # Exact gradient of phi (first term in the chain rule above):
> jax.grad(phi)(theta_opt)
\end{lstlisting}

See \texttt{rgiordan/vittles} on \texttt{github}.

\end{frame}



% %%%%%%%%%%%%%%%%%%%%%%%%%%%%%%%%%%%%%%%%%%%%%%%%%%%%%%%%%%%%%%%%%%%%%%%
% %%%%%%%%%%%%%%%%%%%%%%%%%%%%%%%%%%%%%%%%%%%%%%%%%%%%%%%%%%%%%%%%%%%%%%%
% %%%%%%%%%%%%%%%%%%%%%%%%%%%%%%%%%%%%%%%%%%%%%%%%%%%%%%%%%%%%%%%%%%%%%%%

\begin{frame}{How accurate is the approximation?}

\begin{minipage}{0.45\textwidth}
    By conrolling the curvature,
    we can control the error in the linear approximation.
\end{minipage}
%
\begin{minipage}{0.45\textwidth}
\LinearGraph{1}
\end{minipage}

\vspace{1em}
We provide \textbf{finite-sample theory}
\citep{giordano2019swiss}
showing that
$\abs{\thetafun(\thetahat(\w)) - \thetafunlin(\w)} =
O\left(\vnorm{\frac{1}{N}(\w - \onevec)}^2_2\right) =
O\left(\alpha\right)$ as $\alpha \rightarrow 0$.

\vspace{1em}

\hrulefill

\pause
\vspace{1em}
\textbf{But you don't need to rely on the theory to show non-robustness!}

\vspace{1em}
Our method returns which points to drop.  \textbf{Re-running once}
without those points provides an \textbf{exact lower bound} on the worst-case
sensitivity.

% \pause
% \vspace{1em}
% (Let $\wdom$ denote weight vectors dropping $\alphan$ points, and
% $\w^* := \mathrm{argsup}_{\w \in \wdom} \thetafunlin(\w)$.
% Then $\thetafun(\w^*) \le \sup_{\w \in \wdom} \thetafun(\w)$.)
%
\end{frame}


%%%%%%%%%%%%%%%%%%%%%%%%%%%%%%%%%%%%%%%%%%%%%%%%%%%%%%%%%%%%%%%%%%%%%%%%%
%%%%%%%%%%%%%%%%%%%%%%%%%%%%%%%%%%%%%%%%%%%%%%%%%%%%%%%%%%%%%%%%%%%%%%%%%
%%%%%%%%%%%%%%%%%%%%%%%%%%%%%%%%%%%%%%%%%%%%%%%%%%%%%%%%%%%%%%%%%%%%%%%%%


\begin{frame}{Selected related classical ideas}

Our work is most closely related to ``infinitesimal robustness''
\citep{hampel1974influence, hampel1986robustbook}, which is based on the
influence function, the first term in the Taylor series expansion
of functionals of a distribution function.
%
%\vspace{1em}
\begin{itemize}
    \item We diagnose an existing analysis rather than designing
    estimators.
    \item We use autodiff to automate the influence function computation.
    \item We relate sensitivity to the signal-to-noise ratio, improving intuition.
\end{itemize}

\vspace{1em} Our work is also related to the ``global robustness'' concept of a
``breakdown point'' \citep{huber1981robust, huber:1983:notion}.
%
%\vspace{1em}
\begin{itemize}
    \item We consider data dropping and a particular qualitative conclusion,
    rather than arbitrarily large changes.
    \item We use infinitesimal robustness for automaticity and generality,
    but recover non-local robustness properties by controlling the residual.
\end{itemize}

\vspace{1em} Our work shares some formal similarity with outlier detection
\citep{cook1977detection, belsley2005regression}.
%
\begin{itemize}
    \item We consider $O(N)$-sized sets rather than individual points.
    \item We study a specific intervention, not vaguely defined
    ``outliers.''
    \item We consider settings far beyond regression.
\end{itemize}



\end{frame}


%%%%%%%%%%%%%%%%%%%%%%%%%%%%%%%%%%%%%%%%%%%%%%%%%%%%%%%%%%%%%%%%%%%%%%%%%
%%%%%%%%%%%%%%%%%%%%%%%%%%%%%%%%%%%%%%%%%%%%%%%%%%%%%%%%%%%%%%%%%%%%%%%%%
%%%%%%%%%%%%%%%%%%%%%%%%%%%%%%%%%%%%%%%%%%%%%%%%%%%%%%%%%%%%%%%%%%%%%%%%%


\begin{frame}{A simulation}

For $N=\SimNumObs$ data points, compute the OLS estimator from:

\vspace{1em}
\begin{tabularx}{\textwidth}{YYY}
%\begin{tabular}{ccc}
    Regressors  &   Residuals   &   Responses \\
    $x_n \sim \mathcal{N}(0, \sigma_x^2)$   &
    $\varepsilon_n \sim \mathcal{N}(0, \sigma_\varepsilon^2)$   &
    $y_n = \SimTrueTheta x_n + \varepsilon_n$
\end{tabularx}
%

\SimGridNormalGraph{}

\end{frame}



%%%%%%%%%%%%%%%%%%%%%%%%%%%%%%%%%%%%%%%%%%%%%%%%%%%%%%%%%%%%%%%%%%%%%
%%%%%%%%%%%%%%%%%%%%%%%%%%%%%%%%%%%%%%%%%%%%%%%%%%%%%%%%%%%%%%%%%%%%%
%%%%%%%%%%%%%%%%%%%%%%%%%%%%%%%%%%%%%%%%%%%%%%%%%%%%%%%%%%%%%%%%%%%%%


\begin{frame}[t]{Influence function}

The present work is based on the {\em empirical influence function}.
%
Consider:
%
\begin{itemize}
%
\item True, unknown distribution function $\flim(x) = p(X \le x)$
\item Empirical distribution function $\fhat(x) = \meann \ind{\x_n \le x}$
\item A statistical functional $T(F)$.
%
\end{itemize}

\only<2>{
We estimate with $T(\flim)$ with $T(\fhat)$.

Sample means are an example:
%
\begin{align*}
%
T(F) := \int \x \, F(\dee x).
%
\end{align*}
%

Z-estimators are, too:
%
\begin{align*}
%
T(F) := \theta\textrm{  such that  }\int G(\theta, \x) F(\dee x) = 0.
%
\end{align*}
%
}
\only<3->{

Form an (infinite-dimensional) Taylor series expansion at some $F_0$:
%
\begin{align*}
%
T(F) = T(F_0) + T'(F_0) (F - F_0) + \mathrm{residual}.
%
\end{align*}
%
When the derivative operator takes the form of an integral
%
\begin{align*}
%
T'(F_0) \Delta = \int \infl(\x; F_0) \Delta(\dee x)
%
\end{align*}
%
then $\infl(\x; F_0)$ is known as the {\em influence function}.

{
Where to form the expansion? There are at least two reasonable choices:
%
\begin{itemize}
%
\item The limiting influence function $\infl(\x, \flim)$
\item The empirical influence function $\infl(\x, \fhat)$
%
\end{itemize}
%
}

}


\end{frame}



%%%%%%%%%%%%%%%%%%%%%%%%%%%%%%%%%%%%%%%%%%%%%%%%%%%%%%%%%%%%%%%%%%%%%
%%%%%%%%%%%%%%%%%%%%%%%%%%%%%%%%%%%%%%%%%%%%%%%%%%%%%%%%%%%%%%%%%%%%%
%%%%%%%%%%%%%%%%%%%%%%%%%%%%%%%%%%%%%%%%%%%%%%%%%%%%%%%%%%%%%%%%%%%%%


\begin{frame}[t]{Influence function}

\begin{itemize}
%
\item The limiting influence function (LIF) $\infl(\x, \flim)$
    \begin{itemize}
        \item Used in a lot of classical statistics
            \citep{
                mises1947asymptotic,
                huber1981robust, hampel1986robustbook,
                bickel1993semiparametric}
        \item Unobserved, asymptotic
        \item Requires careful functional analysis
            \citep{reeds1976thesis}
    \end{itemize}
\item The empirical influence function (EIF) $\infl(\x, \fhat)$
    \begin{itemize}
        \item The basis of the present work
            (also \citep{giordano2019swiss, giordano2019higherorder})
        \item Computable, finite-sample
        \item Requires only finite-dimensional calculus
    \end{itemize}
%
\end{itemize}

\vspace{1em}
Typically the {\em semantics} of the EIF derive from study of the LIF.

Example: $\meann (N \infl_n)^2 \approx \var{}{\sqrt{N}\thetafun(\thetahat)}$.

\vspace{1em}
But the EIF measures what happens when you perturb the data at hand.

\vspace{1em}
Other data perturbations will admit an analysis similar to ours!

\end{frame}



%%%%%%%%%%%%%%%%%%%%%%%%%%%%%%%%%%%%%%%%%%%%%%%%%%%%%%%%%%%%%%%%%%%%%
%%%%%%%%%%%%%%%%%%%%%%%%%%%%%%%%%%%%%%%%%%%%%%%%%%%%%%%%%%%%%%%%%%%%%
%%%%%%%%%%%%%%%%%%%%%%%%%%%%%%%%%%%%%%%%%%%%%%%%%%%%%%%%%%%%%%%%%%%%%


\begin{frame}{Local robustness}

The present work is an application of {\em local robustness}.  Consider:
%
\begin{itemize}
%
\item Model parameter $\lambda$  (e.g., data weights $\lambda = \w$)
\item Set of plausible models $\mathcal{S}_\lambda$
    (e.g. $\mathcal{S}_\lambda = W_\alpha$)
\item Estimator $\thetahat(x, \lambda)$ for data $\x$ and
    $\lambda \in \mathcal{S}_\lambda$
    (e.g. a Z-estimator)
%
\end{itemize}
%
% Specify data $x$, a set of plausible models $\lambda \in \mathcal{S}_\lambda$,
% and an estimator $\thetahat(x, \lambda)$.

\hrulefill

\begin{tabular}{ccc}
    Global robustness:   &
%
$
\left(
\inf_{\lambda \in \mathcal{S}_\lambda} \thetahat(x, \lambda),
\sup_{\lambda \in \mathcal{S}_\lambda} \thetahat(x, \lambda)
\right)
$
& (Hard in general!)
\end{tabular}


\hrulefill

\begin{tabular}{cc}
Local robustness:
&
$
\left(
\inf_{\lambda \in \mathcal{S}_\lambda} \thetahat^{lin}(x, \lambda),
\sup_{\lambda \in \mathcal{S}_\lambda} \thetahat^{lin}(x, \lambda)
\right)
$
%
\end{tabular}

...where $\thetahat^{lin}(\x, \lambda) :=
\thetahat^{lin}(\x, \lambda_0) +
    \fracat{\partial \thetahat^{lin}(\x, \lambda)}{\partial \lambda}{\lambda_0}
        (\lambda  - \lambda_0)$.

\hrulefill

\textbf{Many variants are possible!}

\begin{itemize}
    \item Cross-validation \citep{giordano2019swiss}
    \item Prior sensitivity in Bayesian nonparametrics \citep{giordano2021bnp}
    \item Model sensitivity of MCMC output \citep{giordano2018covariances}
    \item Frequentist variances of MCMC posteriors (in progress)
\end{itemize}

% My work emphasizes \textbf{efficient computation of the derivative}
% and \textbf{accuracy as an approximation to global robustness}.

\end{frame}
